\documentclass[11pt]{article}
\usepackage[a4paper,margin=27mm]{geometry}
\usepackage[T1]{fontenc}
\usepackage{lmodern,amsmath,amssymb,amsthm,mathtools,microtype}
\usepackage[hidelinks]{hyperref}
\usepackage{enumitem}
\setlist{nosep}
\setlength{\emergencystretch}{2em}
\newtheorem{theorem}{Theorem}[section]
\newtheorem{lemma}[theorem]{Lemma}
\newtheorem{proposition}[theorem]{Proposition}
\newtheorem{corollary}[theorem]{Corollary}
\theoremstyle{remark}
\newtheorem{remark}[theorem]{Remark}
\newcommand{\kk}{\mathbb K}
\newcommand{\Sym}{\operatorname{Sym}}
\newcommand{\Span}{\operatorname{span}}
\newcommand{\rank}{\operatorname{rank}}
\newcommand{\im}{\operatorname{im}}
\newcommand{\row}{\operatorname{row}}
\newcommand{\perm}{\operatorname{perm}}
\newcommand{\ChowRank}{\operatorname{ChowRank}}
\newcommand{\id}{\operatorname{id}}
\newcommand{\coker}{\operatorname{coker}}
\newcommand{\dd}{\partial}
\title{A computation-free proof of the Chow rank\ of the four-by-four permanent}
\author{Rongyu Xu}
\date{Research revision, October 19, 2026}
\hypersetup{pdftitle={A computation-free proof of the Chow rank of the four-by-four permanent},pdfauthor={Rongyu Xu}}
\begin{document}
\maketitle
\begin{abstract}
We give a self-contained proof that the ordinary Chow rank of the generic
four-by-four permanent is eight in characteristic zero. The two base Koszul
ranks are obtained by classifying cubic prolongations. The uniform quadratic
quotient-kernel lemma is proved first at a coordinate vector by a
missing-derivative argument, then at every nonzero vector by an explicit
row-column scaling and specialization of matrix rank. This replaces the
large-minor computation in the earlier proof. No matrix-rank computation,
finite enumeration, or computer-generated certificate is a premise of the
argument. Optional exact regression tests accompany this research revision.
\end{abstract}

\paragraph{Status and provenance.}
This is an AI-assisted research revision of the author's existing perm4
proof package \cite{original}. Its mathematical argument has received an
internal review, not external peer review or formal proof-assistant
certification. The archived manuscript and its original computational
certificates are not overwritten. This note makes no assertion about border
Chow rank, or equality of Chow ranks for general matrix size.

\section{Statement and the fixed linear map}
Initially let $\kk$ be algebraically closed of characteristic zero. Set
\[
 V=\Span_{\kk}\{x_{ij}:0\leq i,j\leq3\},\qquad \dim V=16,
 \qquad f=\perm_4=\sum_{\sigma\in S_4}\prod_{i=0}^3x_{i,\sigma(i)}.
\]
The ordinary Chow rank of $f$ is the least number of products of four linear
forms whose sum is $f$. Scalar coefficients can be absorbed in one factor.
\begin{theorem}\label{thm:main}
Over every field of characteristic zero,
\[
 \ChowRank(\perm_4)=8.
\]
\end{theorem}

For a quartic $h$, put
\[
 C_h:\Sym^2V^*\longrightarrow\Sym^2V,\qquad D\longmapsto D(h),
 \qquad E_h=\im C_h.
\]
Define the fixed map
\begin{align*}
 \delta:\Sym^2V\otimes V&\longrightarrow V\otimes\Lambda^2 V,\\
 (uv)\otimes w&\longmapsto v\otimes(u\wedge w)+u\otimes(v\wedge w),
\end{align*}
and $K(h)=\delta\circ(C_h\otimes\id_V)$. In particular, $K$ is linear in
$h$, its source and target do not depend on $h$, and
\begin{equation}\label{eq:image-row}
 \im K(h)=\delta(E_h\otimes V),\qquad
 \row K(h)\subseteq E_h\otimes V^*.
\end{equation}
For the second inclusion use the natural apolar pairings: for
$D,D'\in\Sym^2V^*$,
$\langle D',C_h(D)\rangle=D'D(h)=DD'(h)=\langle D,C_h(D')\rangle$.
Thus $C_h$ is self-adjoint under these identifications, and the image of
$K(h)^*$ is contained in $E_h\otimes V^*$.
No unnormalized monomial-basis symmetry is assumed.

\section{The Koszul kernel and the two base ranks}
For $F\subseteq\Sym^2V$ define
\[
 F^{(1)}=\{g\in\Sym^3V:\dd_a g\in F\text{ for every coordinate }a\}.
\]
Here a single index $a$ can stand for a pair $(i,j)$.
\begin{lemma}\label{lem:gradient}
The map $\iota(g)=\sum_a\dd_a g\otimes x_a$ identifies $F^{(1)}$ with
$\ker(\delta|_{F\otimes V})$. Consequently
\[
 \rank\delta(F\otimes V)=16\dim F-\dim F^{(1)}.
\]
\end{lemma}
\begin{proof}
For $\eta=\sum_a q_a\otimes x_a$, the coefficient of $x_b\wedge x_a$
in $\delta\eta$ is $\dd_bq_a-\dd_aq_b$, an element of $V$. Thus
$\delta\eta=0$ precisely when all these mixed-partial equalities hold.
In this case put $g=\frac13\sum_a x_aq_a$. Euler's identity for the
quadratic $q_b$ gives
\[
 \dd_b g=\tfrac13\bigl(q_b+\sum_a x_a\dd_bq_a\bigr)
 =\tfrac13\bigl(q_b+\sum_a x_a\dd_aq_b\bigr)=q_b.
\]
Conversely a gradient has zero image under $\delta$, by equality of mixed
partials. The gradient map on homogeneous cubics is injective in
characteristic zero. Restriction to $F\otimes V$ proves the assertion.
\end{proof}

For the permanent let $E=E_f$. Direct differentiation gives
\begin{equation}\label{eq:E}
 E=\Span\{x_{ia}x_{jb}+x_{ib}x_{ja}: i<j,\ a<b\},\qquad \dim E=36.
\end{equation}
Each generator has its own row-pair and column-pair support. In particular,
a quadratic belongs to $E$ precisely when its coefficients on squares and
same-row or same-column pairs vanish, and its two coefficients on each
rectangle are equal.
Let
\[
 D=\Span\{\perm(X_{I,J}): |I|=|J|=3\}.
\]
\begin{lemma}\label{lem:full-prolongation}
$E^{(1)}=D$ and $\dim D=16$.
\end{lemma}
\begin{proof}
If a cubic monomial has two factors in a common row or a common column,
including a repeated factor, differentiation by the remaining factor
produces a forbidden quadratic monomial. For a fixed derivative and a fixed
remaining monomial there is only one source monomial; its coefficient
therefore cannot be cancelled by another source. The derivative multiplicity
is nonzero in characteristic zero. Such a cubic has zero coefficient.

The surviving monomials are three-edge matchings. Fix their three rows $I$
and three columns $J$. There are six matching monomials. Swapping the columns
assigned to two rows leaves a common third factor. Differentiating by that
factor and applying the rectangle coefficient condition in \eqref{eq:E}
shows that the two cubic coefficients agree. Transpositions connect all six
permutations, so all six coefficients agree. Their sum is
$\perm(X_{I,J})$. Conversely every first derivative of this cubic belongs
to $E$. Distinct pairs $(I,J)$ have disjoint monomial supports, and their
number is $\binom43^2=16$.
\end{proof}

\begin{proposition}\label{prop:base}
$\rank K(f)=560$. For every Chow term $T=\ell_1\ell_2\ell_3\ell_4$,
$\rank K(T)\leq92$, with equality when the factors are independent.
\end{proposition}
\begin{proof}
Lemmas~\ref{lem:gradient} and \ref{lem:full-prolongation} give
$\rank K(f)=16\cdot36-16=560$.
For independent factors change coordinates to $T=y_0y_1y_2y_3$.
Then $E_T$ is the span of the six products $y_iy_j$, $i<j<4$.
Its first prolongation consists of the four squarefree cubics in these four
variables. Indeed a repeated factor or a variable outside their span is
ruled out by differentiating to a forbidden quadratic; all four indicated
cubics qualify. Hence $\rank K(T)=16\cdot6-4=92$.
The entries of the fixed matrix $K(\ell_1\ell_2\ell_3\ell_4)$ are
polynomial in the factor coefficients. Its $93$-minors vanish on the dense
open set of independent factor tuples, so they vanish identically. This
proves the upper bound for dependent and repeated factors as well.
\end{proof}

\section{The missing-derivative lemma}
Write $z=x_{00}$ and set
\[
 P_z=\{g\in\Sym^3V:\dd_{ij}g\in E\text{ for every }(i,j)\ne(0,0)\}.
\]
Unlike $E^{(1)}$, this space imposes no condition on $\dd_{00}g$.
\begin{lemma}\label{lem:missing}
\[
 P_z=D\oplus\kk z^3.
\]
\end{lemma}
\begin{proof}
We first rule out all nonmatching cubic monomials except $z^3$.
A nonmatching monomial with no factor $z$ is excluded just as in
Lemma~\ref{lem:full-prolongation}, since the derivative used is not $\dd_z$.
A monomial $z^2u$ with $u\ne z$ is excluded by $\dd_u$, which produces
$z^2$. For a monomial $zuw$, with $u,w\ne z$, if either $u$ or $w$ shares
a row or column with $z$, differentiating by the other factor leaves a
forbidden pair. This also applies when $u=w$.

It remains to consider $zuw$ where $u$ and $w$ are interior variables
(their row and column indices are nonzero) and share a row or column with
each other. Write $w=x_{ab}$ with $a,b\ne0$. In $\dd_u g$ the
coefficient of $zw$ is a nonzero integer multiple of the coefficient of
$zuw$. Membership in $E$ forces that coefficient to equal the coefficient
of $x_{0b}x_{a0}$. The unique source for the latter coefficient is
\[
 m'=u\,x_{0b}x_{a0}.
\]
If $u,w$ share row $a$, differentiating $m'$ by $x_{0b}\ne z$ leaves
the forbidden same-row pair $u x_{a0}$. If they share column $b$,
differentiating by $x_{a0}\ne z$ leaves the forbidden same-column pair
$u x_{0b}$. Thus $m'$ has coefficient zero, forcing the coefficient of
$zuw$ to be zero too. Repeated interior factors are included in this
argument; their derivative multiplicity is two, not zero.

All remaining monomials are matchings or $z^3$. For fixed three-row and
three-column sets the usual transposition coefficient equalities still
hold whenever their common third factor is not $z$. If $z$ can occur on
this support, exactly one of the nine transposition edges is missing:
the edge between the two matchings using $z$. The six-vertex transposition
graph is $K_{3,3}$, which stays connected after deleting one edge.
Equivalently, the missing equality is supplied by the following three
allowed swaps (where $i,k,a,b\ne0$, $i\ne k$, $a\ne b$):
\begin{align*}
 z x_{ia}x_{kb}
 &\longleftrightarrow x_{0a}x_{i0}x_{kb}\\
 &\longleftrightarrow x_{0a}x_{ib}x_{k0}
 \longleftrightarrow z x_{ib}x_{ka}.
\end{align*}
The common factors in these three comparisons are, respectively,
$x_{kb}$, $x_{0a}$ and $x_{ib}$, all different from $z$. Thus the six
coefficients still agree. Each support contributes a subpermanent, and
the only extra freedom is the coefficient of $z^3$.
The reverse inclusion is immediate, and $z^3$ is not in $D$.
\end{proof}

\section{The uniform quotient-kernel lemma, without a minor}
Let $I=\delta(E\otimes V)$ and define
\[
 Q=\Sym^2V/E,\qquad W=(V\otimes\Lambda^2V)/I,
\]
\[
 \psi_v:Q\longrightarrow W,\qquad [q]\longmapsto[\delta(q\otimes v)].
\]
The map is well defined, and its matrix in fixed bases is linear in $v$.
Since $\dim Q=136-36=100$ and $\delta(v^2\otimes v)=0$, its rank is at
most $99$ for $v\ne0$. Here $[v^2]\ne0$: a nonzero square has at least
one nonzero square coefficient, whereas every element of $E$ has none.

\begin{proposition}\label{prop:uniform}
For every nonzero $v\in V$,
\[
 \ker\psi_v=\kk[v^2],\qquad \rank\psi_v=99.
\]
\end{proposition}
\begin{proof}
First let $v=z$. If $[q]\in\ker\psi_z$, choose
$\eta=\sum_a e_a\otimes x_a\in E\otimes V$ such that
$\delta(q\otimes z-\eta)=0$. Lemma~\ref{lem:gradient} gives a cubic $g$
with
\[
 \dd_a g=-e_a\quad(a\ne z),\qquad \dd_zg=q-e_z.
\]
By Lemma~\ref{lem:missing}, $g=d+cz^3$ for $d\in D$. Since
$\dd_z d\in E$, it follows that $[q]=3c[z^2]$. The reverse inclusion
was already observed, so $\rank\psi_z=99$.

Now take $v=\sum v_{ij}x_{ij}\ne0$. Row and column permutations preserve
$E$ and commute with the construction of $\delta$, so assume $v_{00}\ne0$.
For $t\ne0$ let $\tau_t$ be the row-column scaling
\begin{equation}\label{eq:scaling}
 \tau_t(x_{ij})=t^{\mathbf1_{i\ne0}+\mathbf1_{j\ne0}}x_{ij}.
\end{equation}
Both monomials of every rectangle in \eqref{eq:E} acquire the same weight,
so $\tau_t(E)=E$. Equivariance of $\delta$ then implies
\[
 \psi_{\tau_tv}=\overline\tau_{t,W}\,\psi_v\,
                 \overline\tau_{t,Q}^{-1}\qquad(t\ne0).
\]
Thus $\rank\psi_{\tau_tv}=\rank\psi_v$ for every nonzero $t$.
On the other hand,
\[
 \tau_tv=v_{00}z+t\sum_{j>0}v_{0j}x_{0j}
                  +t\sum_{i>0}v_{i0}x_{i0}
                  +t^2\sum_{i,j>0}v_{ij}x_{ij}
\]
is a polynomial family with specialization $v_{00}z$ at $t=0$.
If its constant rank for $t\ne0$ is $r$, all $(r+1)$-minors vanish for
infinitely many $t$, hence identically, and in particular at $t=0$.
Therefore
\[
 99=\rank\psi_{v_{00}z}\leq\rank\psi_v\leq99.
\]
This proves the assertion. Notice the direction: rank can drop, not rise,
under specialization. No degeneration of a putative kernel vector is used.
\end{proof}

\begin{remark}
The proof does not evaluate, or assert a new evaluation of, the historical
$659\times659$ determinant. It proves the rank statement that determinant
was used for, rendering that calculation unnecessary. Only the finite
six-matching coefficient argument above is used; it is proved explicitly,
not supplied by an enumerator or a rank program.
\end{remark}

\section{Low-rank quadratic structure}
Identify $V=A\otimes B$ using the matrix rows and columns. For a quadratic
$q$ let $M(q)$ be its Hessian, so that its image is the essential variable
space of $q$. In row blocks, a quadratic in $E$ has a matrix
\[
 M(q)=(D_{ij})_{0\leq i,j\leq3},\qquad D_{ii}=0,\quad D_{ij}=D_{ji},
\]
where each $D_{ij}$ is symmetric with zero diagonal.

\begin{lemma}\label{lem:zero-diagonal}
A nonzero symmetric zero-diagonal matrix has rank at least two.
If such a matrix $D$ has rank two, every symmetric zero-diagonal matrix
$H$ with $\im H\subseteq\im D$ is a scalar multiple of $D$.
\end{lemma}
\begin{proof}
A symmetric rank-one matrix is $caa^{\mathsf T}$; its zero diagonal forces
$a=0$, a contradiction. Over the algebraically closed field, a rank-two
symmetric form can be written
$D=ab^{\mathsf T}+ba^{\mathsf T}$ with $a,b$ independent.
Zero diagonal says $a_ib_i=0$, so their nonempty supports are disjoint.
Symmetry and the image condition give
$H=[a\ b]R[a\ b]^{\mathsf T}$ for a symmetric two-by-two matrix $R$.
The diagonal conditions on the two supports force $R_{11}=R_{22}=0$.
Only the off-diagonal coefficient remains.
\end{proof}

\begin{proposition}\label{prop:lowrank}
If $0\ne q\in E$ and $\rank M(q)\leq4$, then
\[
 M(q)=U\otimes D,
\]
where $U,D$ are symmetric zero-diagonal four-by-four matrices of rank two.
In particular, every nonzero element of $E$ has rank at least four.
\end{proposition}
\begin{proof}
Choose a nonzero block $D=D_{ij}$. Its two-row principal block is
\[
 N=\begin{bmatrix}0&D\\D&0\end{bmatrix},\qquad
 2\rank D=\rank N\leq\rank M(q)\leq4.
\]
Lemma~\ref{lem:zero-diagonal} forces $\rank D=2$ and
$\rank N=\rank M(q)=4$. The connecting block columns must lie in
$\im N$: otherwise the rank of the full matrix would exceed that of this
principal block. Hence
$\im D_{ik},\im D_{jk}\subseteq\im D$ for every $k$.
By Lemma~\ref{lem:zero-diagonal}, write
$D_{ik}=a_kD$ and $D_{jk}=b_kD$.

For completeness, if a symmetric matrix is written
$\left[\begin{smallmatrix}N&B\\B^{\mathsf T}&C\end{smallmatrix}\right]$
and has the same rank as $N$, then $B=NX$ for some $X$, and congruence
elimination yields $C=X^{\mathsf T}NX$. Apply this with a generalized
inverse of $D$ satisfying $DGD=D$. The remaining blocks are
\[
 D_{kl}=(a_kb_l+b_ka_l)D.
\]
Thus all blocks are multiples of $D$, giving $M(q)=U\otimes D$.
The block symmetries and zero diagonal give the same properties for $U$.
Finally $4=\rank M(q)=\rank U\,\rank D$, so $\rank U=2$.
The minimum-rank assertion also follows directly from the principal block
$N$ for any nonzero $q$, without the hypothesis $\rank M(q)\leq4$.
\end{proof}

\begin{corollary}\label{cor:intersection}
For $L\subseteq V$ with $\dim L\leq4$,
\[
 \dim(E\cap\Sym^2L)\leq1.
\]
In particular, for every Chow term $T$, $\dim(E\cap E_T)\leq1$.
\end{corollary}
\begin{proof}
If the intersection is nonzero, Proposition~\ref{prop:lowrank} implies
$\dim L=4$ and
$L=\im U\otimes\im D$ for any nonzero $q$ in the intersection.
The two tensor factors of a nonzero product subspace are determined by
that subspace: contractions against all linear functionals on the other
factor recover each one. Consequently every nonzero $q$ has the same
two-dimensional image spaces for $U$ and $D$. By
Lemma~\ref{lem:zero-diagonal}, each of these matrices is unique up to
scale inside its fixed image. Thus all the quadratics are proportional.
For a Chow term, $E_T\subseteq\Sym^2\Span(\ell_1,\ldots,\ell_4)$,
so the last assertion follows, including dependent factors.
\end{proof}

\section{Adding a quadratic direction and subtracting a Chow term}
\begin{lemma}\label{lem:extension}
For every $q\in\Sym^2V\setminus E$,
\[
 \dim\delta((E+\kk q)\otimes V)\geq575.
\]
\end{lemma}
\begin{proof}
Let $\theta_q:V\to W$ send $v$ to $[\delta(q\otimes v)]$.
If independent vectors $v,w$ belonged to its kernel,
Proposition~\ref{prop:uniform} would give
$[q]=a[v^2]=b[w^2]$ for nonzero scalars $a,b$.
Then $0\ne av^2-bw^2\in E$ has rank at most two, contrary to
Proposition~\ref{prop:lowrank}. Hence $\dim\ker\theta_q\leq1$ and
$\rank\theta_q\geq15$. Its image is the new image modulo $I$, so the
claimed dimension is at least $560+15=575$.
\end{proof}

\begin{lemma}\label{lem:rank-inequality}
For matrices $A,B$ with the same source and target,
\[
 \rank(A-B)\geq\rank[\,A\ B\,]
 +\rank\begin{bmatrix}A\\B\end{bmatrix}-\rank A-\rank B.
\]
\end{lemma}
\begin{proof}
Write $a=\rank A$, $b=\rank B$, and
$c=\dim(\im A\cap\im B)$, $d=\dim(\row A\cap\row B)$.
The concatenation ranks are $a+b-c$ and $a+b-d$.
The map $x\mapsto Ax$ sends $\ker(A-B)$ into $\im A\cap\im B$
with kernel $\ker A\cap\ker B$. If the common source has dimension $n$,
this last kernel has dimension $n-(a+b-d)$.
Thus $\dim\ker(A-B)\leq c+n-a-b+d$, giving the displayed inequality.
\end{proof}

\begin{proposition}\label{prop:subtract}
For every nonzero Chow term $T$,
\[
 \rank K(f-T)\geq559.
\]
\end{proposition}
\begin{proof}
Put $A=K(f)$, $B=K(T)$ and $F=E_T$.
Corollary~\ref{cor:intersection} leaves two cases.
If $E\cap F=0$, \eqref{eq:image-row} makes the row spaces of $A$ and
$B$ disjoint. Therefore the vertical concatenation rank equals
$\rank A+\rank B$, and Lemma~\ref{lem:rank-inequality} gives
\[
 \rank(A-B)\geq\rank[\,A\ B\,]\geq560.
\]
If $\dim(E\cap F)=1$, a nonzero common quadratic has rank four.
The factors of $T$ consequently span four dimensions, so they are
independent. Hence $\dim F=6$ and $\rank B=92$.
Choose $q\in F\setminus E$. Lemma~\ref{lem:extension} implies
\[
 \rank[\,A\ B\,]\geq575.
\]
The intersection of the two row spaces is contained in
$(E\cap F)\otimes V^*$, of dimension $16$. Thus
\[
 \rank\begin{bmatrix}A\\B\end{bmatrix}\geq560+92-16=636.
\]
Lemma~\ref{lem:rank-inequality} now yields
\[
 \rank(A-B)\geq575+636-560-92=559.
\]
Since $K$ is linear in its quartic input, $A-B=K(f-T)$.
\end{proof}

\begin{proof}[Lower bound in Theorem~\ref{thm:main}]
If $f$ were a sum of at most seven nonzero Chow terms, subtract one term.
The remainder would be a sum of at most six Chow terms and would have
Koszul rank at most $6\cdot92=552$, by linearity and subadditivity of
matrix rank. Proposition~\ref{prop:subtract} instead gives a lower bound
$559>552$. This contradiction proves $\ChowRank(f)\geq8$.
\end{proof}

\section{The eight-term upper bound and the logical boundary}
Glynn's identity \cite{glynn} in size four is
\begin{equation}\label{eq:glynn}
 \perm_4=\frac18\sum_{\substack{\epsilon\in\{1,-1\}^4\\\epsilon_0=1}}
       \left(\prod_{i=0}^3\epsilon_i\right)
       \prod_{j=0}^3\left(\sum_{i=0}^3\epsilon_i x_{ij}\right).
\end{equation}
Here is a direct verification, independent of an expansion program.
Every monomial in the right side chooses one row in each column.
Let $m_i$ be the number of choices of row $i$. Its coefficient is
\[
 \frac18\prod_{i=1}^3\left(\sum_{\epsilon_i=\pm1}
                                 \epsilon_i^{m_i+1}\right).
\]
It vanishes unless $m_1,m_2,m_3$ are odd. Since the $m_i$ are nonnegative
and sum to four, this forces $m_1=m_2=m_3=1$ and $m_0=1$.
Exactly the permutation monomials remain, each with coefficient one.
There are eight summands in \eqref{eq:glynn}, each a Chow term. This proves
the upper bound and completes Theorem~\ref{thm:main} over algebraically
closed fields of characteristic zero. For any characteristic-zero field,
extend scalars to its algebraic closure for the lower bound. The upper
bound is defined over $\mathbb Q$, so the equality holds over that field too.

\paragraph{What has been replaced.}
The old ranks $560$ and $92$ are now proved by
Lemmas~\ref{lem:gradient} and \ref{lem:full-prolongation} and
Proposition~\ref{prop:base}. The rank-$99$ calculation is replaced by
Lemma~\ref{lem:missing} and the scaling \eqref{eq:scaling}. The low-rank
classification, quadratic extension, and final subtraction argument are
included in full so the proof is self-contained.
Optional scripts in this revision check coefficient rules and ranks over
$\mathbb Q$ and two prime fields. Their outputs are not invoked by any
lemma, proposition, or theorem above.

\paragraph{Why not substitute a tensor-rank bound.}
The recursive Koszul result \cite{hanjukim} concerns the tensor rank of
$\perm_4$, a different rank notion. In fact the row-multihomogeneous
component of the single Chow term
$\prod_j(\sum_i x_{ij})$ is $\perm_4$ itself. Thus projection to one
factor from each row does not send an arbitrary Chow term to a simple
tensor. No transfer from that tensor-rank theorem is assumed here.

\begin{thebibliography}{9}
\raggedright
\bibitem{original}
R.~Xu, \emph{The Chow Rank of the Four-by-Four Permanent},
archived manuscript in \texttt{permanent-chow-rank-small-n},
\texttt{perm4/paper/source/n4\_body.tex}, Git blob
\texttt{e983317e31e0a5a6903df0e600cacafcc7190549}.
The research branch before this revision is commit
\texttt{9dd3368a0973ac24a8f63b6ea818a89bda8e1139}.
\bibitem{glynn}
D.~G.~Glynn, \emph{The permanent of a square matrix},
European J. Combin. \textbf{31} (2010), 1887--1891,
doi:10.1016/j.ejc.2010.01.010.
\bibitem{hanjukim}
J.~I.~Han, J.-H.~Ju and Y.~Kim,
\emph{Recursive Koszul flattenings of determinant and permanent tensors},
arXiv:2503.12032 (2025).
\end{thebibliography}
\end{document}
